\section{Core Setups II: Icebergs, Absorption, and Robots}

\subsection{Iceberg-like behavior and hidden absorption}

\subsubsection{What the setup is}

The iceberg setup is one of the most important bridges from naive DOM reading to genuine
microstructure thinking. A visible quote of fixed size keeps reappearing while executed volume
through that level clearly exceeds the displayed amount. In other words, the tape says much
more traded than the book says was visible.

\subsubsection{Why it is believed to work}

The basic hypothesis is that a strong passive participant wants to transact without revealing
full size. That participant is often serious enough to matter because hidden replenishment is
not random retail behavior. If price repeatedly hits the level and the visible slice keeps
refreshing, the trader infers that a large participant is still active.

There are two main paths from this observation:

\begin{itemize}[leftmargin=1.4em]
  \item \term{Fade logic}: the iceberg absorbs flow and price bounces away.
  \item \term{Break logic}: the iceberg is eventually exhausted, and the break after exhaustion can
        be powerful because the market just spent significant energy clearing one serious seller
        or buyer.
\end{itemize}

\subsubsection{What to watch}

\paragraph{Order book.}
Repeated visible replenishment at roughly the same price.

\paragraph{Tape / prints.}
Executed volume far greater than visible size, ideally observed quickly enough that the mismatch
is obvious in real time.

\paragraph{Price behavior.}
Does price stall despite aggressive hitting, or does it keep leaning and finally burst through?
The same iceberg can support either a bounce or a breakout depending on who finally wins.

\subsubsection{Traps and ambiguity}

The transcripts themselves show why this is tricky: sometimes the trader sees an iceberg before
the platform formally marks it; sometimes the apparent iceberg is mixed with additional visible
orders; sometimes the real hidden size is unknown. Therefore the iceberg should not be coded as
``iceberg detected equals trade.'' It should be coded as ``a meaningful hidden-liquidity event
is likely present; now interpret direction with context.''

\subsubsection{Algorithmic translation}

Observable signals could include:

\begin{itemize}[leftmargin=1.4em]
  \item cumulative traded volume at a price divided by concurrent displayed size,
  \item refresh count at the same price within a short window,
  \item number of failed attempts to dislodge the level,
  \item post-absorption displacement away from or through the price.
\end{itemize}

This setup requires prints / executions. Order-book-only historical data is insufficient for a
credible iceberg study.

\subsection{Liquidity absorption as a broader concept}

Icebergs are one special case of absorption. More broadly, absorption means aggressive activity
arrives, but price does not travel as much as it ``should'' if passive resistance were weak.
The practical lesson is the same: flow must be interpreted relative to price response. Heavy
printing by itself is not strength. Heavy printing with no progress can be hidden opposition.

This is a core reason why the future repository must preserve both message streams. Without
trades, the system cannot distinguish between visible size that merely exists and visible or
hidden size that is actually overpowering aggression.

\subsection{Robot-following setups}

\subsubsection{What the setup is}

The robot setup in the transcripts is not a vague claim that ``algos exist.'' It is a very
practical observation that a recurring participant can often be seen acting with machine-like
regularity: placing similar size, repeating the same behavior, pushing incrementally, or
consistently nudging price in one direction.

\subsubsection{Why it is believed to work}

If a participant is algorithmic, persistent, and effective, its footprint can become tradable.
The trader does not need to know the source code of the robot. The trader only needs to infer a
reliable behavior pattern: how aggressive it is, how quickly it refreshes, whether it walks the
book, whether it defends a side, and whether it actually moves price rather than only creating
noise.

\subsubsection{What makes a robot worth following}

The transcripts evaluate a robot by strength, aggression, size, and effectiveness. That is a
useful research checklist:

\begin{itemize}[leftmargin=1.4em]
  \item Does the participant act repeatedly enough to be recognized?
  \item Does its activity correlate with actual price movement?
  \item Is it operating in a liquid instrument where other participants matter, or in a thin
        book where almost anything looks dominant?
  \item Is the trader trading with the robot's pressure or trying to fade it?
\end{itemize}

\subsubsection{Clean versus unclean examples}

Clean robot setups involve a participant whose behavior is legible and effective. Unclean cases
involve projected intentionality: the trader sees a few repeated prints or orders and invents a
robot story that is not supported by persistent evidence.

\subsubsection{Algorithmic translation}

Possible features:

\begin{itemize}[leftmargin=1.4em]
  \item repeated order sizes and refresh cadence at stable intervals,
  \item repeated one-tick nudging patterns,
  \item correlation between one participant-like behavior cluster and subsequent price drift,
  \item side persistence over multiple seconds or minutes,
  \item cancellation / reposting behavior that suggests queue management rather than discretionary
        clicking.
\end{itemize}

This is an area where pattern mining and clustering may later be useful, because the human
notion of ``robot behavior'' often refers to sequence style rather than one isolated snapshot.
